\begin{theorem}[admissible 边界 Gödel 码的乘法型二元线性子码结构]\label{thm:conclusion-dyadic-boundary-godel-admissible-multiplicative-linear-subcode}
\leanverified{paper\_conclusion\_dyadic\_boundary\_godel\_admissible\_multiplicative\_linear\_subcode}
固定 \(m,n\ge 1\)。沿用定理 \ref{thm:conclusion-boundary-godel-syndrome-completeness-linear-decode} 的记号，记
\[
\mathcal G_{m,n}
=
\{\,N\in \mathfrak Q_m^{\partial}:\ \delta_{n-1}(N)=1\,\}.
\]
则 \(\mathcal G_{m,n}\) 是 \(\mathfrak Q_m^{\partial}\) 的一个子群，并且有规范同构
\[
\mathcal G_{m,n}
\cong
\operatorname{im}\partial_n
\cong
C_n(Q_m;\FF_2)
\cong
(\ZZ/2\ZZ)^{2^{mn}}.
\]
特别地，
\[
|\mathcal G_{m,n}|=2^{2^{mn}}.
\]
\end{theorem}

\begin{proof}
由定理 \ref{thm:conclusion-boundary-godel-syndrome-completeness-linear-decode}，
\[
\mathcal G_{m,n}=\delta_{n-1}^{-1}(1).
\]
由于 \(\delta_{n-1}\) 是平方类边界群上的同态，\(\delta_{n-1}^{-1}(1)\) 必为子群。

另一方面，定理 \ref{thm:conclusion-dyadic-boundary-squareclass-syndrome-unique-fill} 表明
\[
N\in \mathcal G_{m,n}
\iff
\mathbf G_{n-1}^{-1}(N)\in \operatorname{im}\partial_n,
\]
且由于顶维边界算子 \(\partial_n\) 单射，\(\operatorname{im}\partial_n\cong C_n(Q_m;\FF_2)\)。
又 \(Q_m\) 的顶维胞腔数恰为 \(2^{mn}\)，故
\[
C_n(Q_m;\FF_2)\cong (\ZZ/2\ZZ)^{2^{mn}}.
\]
由此得到所示全部同构与基数公式。证毕。
\end{proof}

\begin{theorem}[综合征层的等势陪集几何]\label{thm:conclusion-dyadic-boundary-godel-syndrome-fibers-are-equipotent-cosets}
\leanverified{paper\_conclusion\_dyadic\_boundary\_godel\_syndrome\_fibers\_are\_equipotent\_cosets}
在上述记号下，对任意
\[
s\in \operatorname{im}\delta_{n-1},
\]
只要纤维 \(\delta_{n-1}^{-1}(s)\) 非空，它便是 \(\mathcal G_{m,n}\) 的一个陪集。更精确地，若 \(H_0\in\mathfrak Q_m^\partial\) 满足
\[
\delta_{n-1}(H_0)=s,
\]
则
\[
\delta_{n-1}^{-1}(s)=H_0\mathcal G_{m,n}.
\]
因此所有非空综合征层具有完全相同的基数
\[
\bigl|\delta_{n-1}^{-1}(s)\bigr|
=
|\mathcal G_{m,n}|
=
2^{2^{mn}}.
\]
\end{theorem}

\begin{proof}
由于 \(\delta_{n-1}\) 是群同态，对任意 \(H\in\mathfrak Q_m^\partial\) 有
\[
\delta_{n-1}(H)=s
\iff
\delta_{n-1}(H_0^{-1}H)=1
\iff
H_0^{-1}H\in \mathcal G_{m,n}.
\]
这正等价于
\[
H\in H_0\mathcal G_{m,n}.
\]
故非空纤维恰为相应的陪集，其基数与 \(\mathcal G_{m,n}\) 相同。最后由定理 \ref{thm:conclusion-dyadic-boundary-godel-admissible-multiplicative-linear-subcode} 即得显式大小。证毕。
\end{proof}

\begin{corollary}[dyadic 边界 Gödel 码的精确参数与校验秩闭式]\label{cor:conclusion-dyadic-boundary-godel-code-parameters-and-check-rank}
\leanverified{paper\_conclusion\_dyadic\_boundary\_godel\_code\_parameters\_and\_check\_rank}
记
\[
N_{m,n}:=\#F_{n-1}(Q_m),
\qquad
K_{m,n}:=\log_2|\mathcal G_{m,n}|,
\qquad
H_{m,n}:=N_{m,n}-K_{m,n}.
\]
则有闭式
\[
N_{m,n}=n(2^m+1)2^{m(n-1)},
\qquad
K_{m,n}=2^{mn},
\]
以及
\[
H_{m,n}=(n-1)2^{mn}+n2^{m(n-1)}.
\]
因此码率满足
\[
R_{m,n}:=\frac{K_{m,n}}{N_{m,n}}
=
\frac{1}{n(1+2^{-m})}
\longrightarrow
\frac1n
\qquad (m\to\infty).
\]
\end{corollary}

\begin{proof}
长度公式与码率极限由定理 \ref{thm:conclusion-dyadic-boundary-code-dual-realization-parameters} 已给出。又由定理 \ref{thm:conclusion-dyadic-boundary-godel-admissible-multiplicative-linear-subcode}，
\[
K_{m,n}=\log_2|\mathcal G_{m,n}|=2^{mn}.
\]
故
\[
H_{m,n}=N_{m,n}-K_{m,n}
=
n(2^m+1)2^{m(n-1)}-2^{mn}
=
(n-1)2^{mn}+n2^{m(n-1)}.
\]
再代入 \(R_{m,n}=K_{m,n}/N_{m,n}\) 即得所示闭式。证毕。
\end{proof}

\begin{theorem}[dyadic path-consistency 审计的精确查询秩等于校验秩]\label{thm:conclusion-dyadic-boundary-path-consistency-query-rank-equals-check-rank}
\leanverified{paper\_conclusion\_dyadic\_boundary\_path\_consistency\_query\_rank\_equals\_check\_rank}
令 \(\Gamma_{m,n}\) 为 \(Q_m\) 的 \(n\)-胞腔邻接图再附加外部顶点 \(\infty\) 后所得图；其边与全部取向 \((n-1)\)-胞腔一一对应。对任意素数 \(p\)，记
\[
Q_{m,n}^{\mathrm{cyc}}(p)
\]
为在最坏情形下判定模 \(p\) 路径一致性所需的最少回路查询数。则
\[
Q_{m,n}^{\mathrm{cyc}}(p)=H_{m,n}.
\]
\end{theorem}

\begin{proof}
图 \(\Gamma_{m,n}\) 的边数恰为
\[
|E(\Gamma_{m,n})|=N_{m,n},
\]
而顶点数为全部 \(n\)-胞腔外加外部顶点，故
\[
|V(\Gamma_{m,n})|=2^{mn}+1.
\]
由于每个边界 \(n\)-胞腔都与 \(\infty\) 相连，而内部 \(n\)-胞腔通过共 \((n-1)\)-面的邻接连通，\(\Gamma_{m,n}\) 连通。于是
\[
\rank_{\ZZ}H_1(\Gamma_{m,n};\ZZ)
=
|E(\Gamma_{m,n})|-|V(\Gamma_{m,n})|+1
=
N_{m,n}-2^{mn}
=
H_{m,n}.
\]
再由定理 \ref{thm:conclusion-boundary-path-independence-query-complexity-equals-betti}，
\[
Q_{m,n}^{\mathrm{cyc}}(p)=\rank_{\ZZ}H_1(\Gamma_{m,n};\ZZ),
\]
从而得到
\[
Q_{m,n}^{\mathrm{cyc}}(p)=H_{m,n}.
\]
证毕。
\end{proof}

\begin{theorem}[承诺回放与敌手式核验之间的严格复杂度分裂]\label{thm:conclusion-dyadic-boundary-constructive-replay-vs-adversarial-verification-split}
\leanverified{paper\_conclusion\_dyadic\_boundary\_constructive\_replay\_vs\_adversarial\_verification\_split}
在同一 dyadic 边界全息模型中，存在如下严格分裂：
\begin{enumerate}
  \item 若输入已承诺属于 \(\mathcal G_{m,n}\)，则相应 bulk 可在
  \[
  O\bigl(|V(\Gamma_{m,n})|+|E(\Gamma_{m,n})|\bigr)
  =
  O\bigl(2^{mn}+N_{m,n}\bigr)
  \]
  时间内唯一恢复；
  \item 若不预设输入属于 \(\mathcal G_{m,n}\)，而要求仅通过回路查询判定其是否 admissible，则最坏情形必须进行
  \[
  H_{m,n}
  \]
  次查询。
\end{enumerate}
特别地，在固定 \(n\) 而 \(m\to\infty\) 时，承诺回放仍为线性规模，而敌手式核验的最坏复杂度为
\[
\Theta(2^{mn}).
\]
\end{theorem}

\begin{proof}
第一条正是定理 \ref{thm:conclusion-boundary-godel-syndrome-completeness-linear-decode} 的线性时间 bulk 译码。第二条由定理 \ref{thm:conclusion-dyadic-boundary-path-consistency-query-rank-equals-check-rank} 立即给出。再由推论 \ref{cor:conclusion-dyadic-boundary-godel-code-parameters-and-check-rank}，
\[
H_{m,n}=(n-1)2^{mn}+n2^{m(n-1)}=\Theta(2^{mn})
\]
（固定 \(n\)）即得最后一式。证毕。
\end{proof}

\begin{corollary}[可见信息密度与敌手审计密度的精确互补律]\label{cor:conclusion-dyadic-boundary-visible-audit-density-complementarity}
\leanverified{paper\_conclusion\_dyadic\_boundary\_visible\_audit\_density\_complementarity}
定义
\[
\rho_{m,n}^{\mathrm{vis}}:=\frac{K_{m,n}}{N_{m,n}},
\qquad
\rho_{m,n}^{\mathrm{aud}}:=\frac{Q_{m,n}^{\mathrm{cyc}}(2)}{N_{m,n}}.
\]
则有严格恒等式
\[
\rho_{m,n}^{\mathrm{vis}}+\rho_{m,n}^{\mathrm{aud}}=1.
\]
等价地，
\[
\rho_{m,n}^{\mathrm{vis}}=\frac{1}{n(1+2^{-m})},
\qquad
\rho_{m,n}^{\mathrm{aud}}=1-\frac{1}{n(1+2^{-m})}.
\]
因此，当 \(m\to\infty\) 时，
\[
\rho_{m,n}^{\mathrm{vis}}\to \frac1n,
\qquad
\rho_{m,n}^{\mathrm{aud}}\to 1-\frac1n.
\]
\end{corollary}

\begin{proof}
由推论 \ref{cor:conclusion-dyadic-boundary-godel-code-parameters-and-check-rank} 与定理 \ref{thm:conclusion-dyadic-boundary-path-consistency-query-rank-equals-check-rank}，
\[
K_{m,n}+Q_{m,n}^{\mathrm{cyc}}(2)=K_{m,n}+H_{m,n}=N_{m,n}.
\]
两边除以 \(N_{m,n}\) 即得
\[
\rho_{m,n}^{\mathrm{vis}}+\rho_{m,n}^{\mathrm{aud}}=1.
\]
再代入 \(K_{m,n}/N_{m,n}=1/[n(1+2^{-m})]\) 即得显式闭式与极限。证毕。
\end{proof}
